% Chapter — Radio-class sensors: external nav fix and altimeter (FR-23,
% A-6). Follows the mandatory FR-29 six-section template: purpose;
% assumptions and domain of validity; derivation; discretization and
% implementation; validation evidence; references. Written ahead of the
% implementation (Phase 6): this chapter is the binding specification for
% the external-nav-fix and altimeter sensor modules. The implementing
% change adds the chapter-manifest entry and the \testid{} evidence table.
% The equation labels eq:radio:fix, eq:radio:gm, eq:radio:alt, and
% eq:radio:chi2 are intended to be echoed verbatim in the module source
% (FR-29 traceability).
\chapter{Sensors: External Nav Fix and Altimeter}
\label{ch:sensors-radio}

\section{Purpose}
\label{sec:radio:purpose}

This chapter specifies the two radio-class FR-23 sensors. The
\emph{external nav fix} is the project's generalization of GNSS
(assumption A-6): an abstract service that delivers position and
velocity fixes in the GCRF with configurable Gaussian errors, so the
same sensor stands in for GPS, a future LNSS-class lunar service, or
ground-based radiometric tracking without re-architecting. The
\emph{altimeter} delivers a scalar geodetic-altitude measurement over
the central body's reference ellipsoid with a per-run bias and white
noise. Both feed the Chapter~\ref{ch:ekf} reference EKF as aiding
updates, and both are gated by Phase~6 exit criterion~6: error
statistics inside two-sided 95\,\% chi-square bounds over 1{,}000 seeded
draws, with the exact statistic defined in
Section~\ref{sec:radio:stats}.

\section{Assumptions and domain of validity}
\label{sec:radio:domain}

Assumptions:
\begin{enumerate}
  \item \textbf{Fix-level abstraction.} The nav fix models the
    \emph{output} of an external navigation service, not its
    measurements: no pseudoranges, constellation geometry, dilution of
    precision, clock states, or atmospheric delays. Error anisotropy
    that a real service exhibits (e.g.\ vertical inflation) is
    expressible through the per-axis sigmas; error correlation with
    geometry is not, and analyses that depend on it are outside this
    model's domain.
  \item \textbf{GCRF axes.} Fix errors are drawn per GCRF axis. A
    service whose natural error frame is local (east/north/up) is
    approximated by choosing equivalent GCRF sigmas; the difference
    matters only for strongly anisotropic configurations.
  \item \textbf{Optional correlated component.} A first-order
    Gauss--Markov error component exists as an option on both the
    position and the velocity fix (Section~\ref{sec:radio:fix}),
    defaulted off. It models slowly varying service biases
    (ephemeris/clock residuals). When enabled, the chi-square gate of
    Section~\ref{sec:radio:stats} applies only to the white part
    (the gate scenario disables the correlated part); a separate
    deterministic test pins the recursion.
  \item \textbf{Altimeter over the reference ellipsoid.} The altimeter
    measures geodetic altitude over the central body's reference
    ellipsoid exactly as the environment models do
    (Chapter~\ref{ch:environment}): the Bowring two-pass conversion of
    Chapter~\ref{ch:harrispriester},
    equation~\eqref{eq:hp:geodetic}, applied to the Earth-fixed (or
    planet-fixed) truth position via the Chapter~\ref{ch:frames}
    rotation. Terrain, tides, and geoid undulation are not modeled ---
    the instrument is a radar altimeter over the ellipsoid, and any
    scenario needing terrain-relative altitude is out of scope.
  \item \textbf{No latency.} Fixes apply at their sample instant;
    transport and processing latency for realism studies is provided
    generically by the FR-25 \texttt{latency\_cycles} mechanism
    (Chapter~\ref{ch:gnc-builtin}), not duplicated per sensor.
\end{enumerate}

Domain of validity: any state the core can propagate; the altimeter
additionally requires the truth position inside its configured operating
band $h \in [h_{\min}, h_{\max}]$ (outside it the sample is emitted with
$\mathrm{valid} = \mathrm{false}$, the flag semantics of
Chapter~\ref{ch:sensors-optical}) and outside the Bowring conversion's
polar-axis neighborhood only in the sense documented in
Chapter~\ref{ch:harrispriester}. Non-positive sigmas, $\tau_c \le 0$, or
$h_{\min} \ge h_{\max}$ are rejected by the FR-15 validator; the in-loop
model never throws.

\section{Derivation}
\label{sec:radio:derivation}

\subsection{External nav fix}
\label{sec:radio:fix}

At each sample (major-cycle schedule, D-5), with truth position
$\vv{r}^{I}$ and velocity $\vv{v}^{I}$ (the logged \texttt{truth}
channels),
\begin{equation}
  \boxed{\;
  \vv{r}_{\mathrm{meas}} = \vv{r}^{I} + \vv{c}_{r,k} + \vv{\eta}_{r,k},
  \qquad
  \vv{v}_{\mathrm{meas}} = \vv{v}^{I} + \vv{c}_{v,k} + \vv{\eta}_{v,k},
  \;}
  \label{eq:radio:fix}
\end{equation}
with independent white errors per GCRF axis $a \in \{x,y,z\}$,
\begin{equation}
  \eta_{r,k,a} \sim \mathcal{N}(0, \sigma_{r,a}^2),
  \qquad
  \eta_{v,k,a} \sim \mathcal{N}(0, \sigma_{v,a}^2),
  \label{eq:radio:white}
\end{equation}
and the optional first-order Gauss--Markov components (defaulted off,
$\vv{c} \equiv \vv{0}$) following exactly the recursion of
Chapter~\ref{ch:sensors-imu}, equation~\eqref{eq:imu:gm}:
\begin{equation}
  \vv{c}_{k} = \varphi\,\vv{c}_{k-1} + \vv{w}_k,
  \qquad
  \varphi = e^{-\Delta t_s/\tau_c},
  \qquad
  \vv{w}_k \sim \mathcal{N}\bigl(\vv{0},\,
    \sigma_{\mathrm{GM}}^2(1 - \varphi^2)\,\mat{I}_3\bigr),
  \qquad
  \vv{c}_{0} \sim \mathcal{N}\bigl(\vv{0},\,
    \sigma_{\mathrm{GM}}^2\,\mat{I}_3\bigr),
  \label{eq:radio:gm}
\end{equation}
with separate $(\sigma_{\mathrm{GM}}, \tau_c)$ pairs for the position
and velocity components. The exactness of the discretization and its
stationary initialization are established in
Chapter~\ref{ch:sensors-imu}. The fix carries
$\mathrm{valid} = \mathrm{true}$ unconditionally in v1; outage modeling
is a configuration concern deferred until a use case demands it (PRD
Non-Goals discipline), and the flag exists in the interface so adding it
later breaks nothing.

\subsection{Altimeter}
\label{sec:radio:alt}

With $h(\vv{r})$ the geodetic altitude of the truth position over the
central body's reference ellipsoid --- the planet-fixed position
$\dcm{I}{E}\,\vv{r}^{I}$ (Chapter~\ref{ch:frames}) put through the
Bowring two-pass conversion of
equation~\eqref{eq:hp:geodetic} with the body's $(a, 1/f)$; for a
spherical body ($f = 0$) the conversion degenerates exactly to
$h = \norm{\vv{r}} - a$ --- the measurement is
\begin{equation}
  \boxed{\;
  h_{\mathrm{meas},k} = h(\vv{r}^{I}_k) + b_h + \eta_{h,k},
  \qquad
  b_h \sim \mathcal{N}(0, \sigma_{b}^2) \text{ once per run},
  \qquad
  \eta_{h,k} \sim \mathcal{N}(0, \sigma_{h}^2),
  \;}
  \label{eq:radio:alt}
\end{equation}
with validity
\begin{equation}
  \mathrm{valid}_k = \bigl[h_{\min} \le h(\vv{r}^{I}_k) \le
    h_{\max}\bigr],
  \label{eq:radio:altgate}
\end{equation}
the configured operating band of the instrument. The measurement is
computed and logged regardless of the flag, and draws are unconditional,
per the sensor-wide semantics of Chapter~\ref{ch:sensors-optical}.

\subsection{Acceptance statistics (exit criterion 6)}
\label{sec:radio:stats}

The gate scenario holds truth fixed, disables the correlated components
($\sigma_{\mathrm{GM}} = 0$ for the fix, $\sigma_b = 0$ for the
altimeter), and collects $M = 1000$ seeded samples. Per sample $k$ the
statistics are the exact normalized quadratic forms
\begin{equation}
  q^{(k)}_{r} = \sum_{a}
    \frac{(r_{\mathrm{meas},a} - r^{I}_a)^2_{(k)}}{\sigma_{r,a}^2}
    \sim \chi^2_3,
  \qquad
  q^{(k)}_{v} = \sum_{a}
    \frac{(v_{\mathrm{meas},a} - v^{I}_a)^2_{(k)}}{\sigma_{v,a}^2}
    \sim \chi^2_3,
  \qquad
  q^{(k)}_{h} =
    \frac{(h_{\mathrm{meas}} - h(\vv{r}^{I}))^2_{(k)}}{\sigma_{h}^2}
    \sim \chi^2_1,
  \label{eq:radio:chi2}
\end{equation}
each i.i.d.\ across samples by construction, and the gates are two-sided
on the ensemble means:
\begin{equation}
  \frac{\chi^2_{0.025}(nM)}{M} \le \frac{1}{M} \sum_{k=1}^{M} q^{(k)}
    \le \frac{\chi^2_{0.975}(nM)}{M},
  \qquad
  \begin{cases}
    [2.850,\; 3.154] & n = 3 \text{ (fix position, fix velocity)},\\[2pt]
    [0.914,\; 1.090] & n = 1 \text{ (altimeter)},
  \end{cases}
  \label{eq:radio:bounds}
\end{equation}
for $M = 1000$, quantiles evaluated by the Wilson--Hilferty form of
Chapter~\ref{ch:ekf}, equation~\eqref{eq:ekf:wh}. These are exactly the
statistics the committed pytest driver computes from the logged
\texttt{sensors.*} and \texttt{truth} channels; the bias and
Gauss--Markov terms are pinned separately by deterministic tests (the
seeded draw values themselves), because a per-run bias makes the
per-sample statistics of equation~\eqref{eq:radio:chi2} dependent and
the mean gate invalid.

\section{Discretization and implementation}
\label{sec:radio:implementation}

Binding implementation notes for the Phase~6 modules:
\begin{enumerate}
  \item \textbf{Modules and traceability.} Both sensors are core-side
    \texttt{ISensor} modules on the named streams
    \texttt{sensors.navfix} and \texttt{sensors.alt} (D-9). Source
    comments echo \texttt{eq:radio:fix}, \texttt{eq:radio:gm},
    \texttt{eq:radio:alt}, and \texttt{eq:radio:chi2} verbatim.
  \item \textbf{Draw schedule, normative.} Nav fix, per sample and in
    order: position Gauss--Markov drive $(x,y,z)$ if enabled, velocity
    Gauss--Markov drive $(x,y,z)$ if enabled, position white $(x,y,z)$,
    velocity white $(x,y,z)$. Enabling a Gauss--Markov component also
    consumes its stationary-initialization draws at run start (position
    then velocity). Altimeter: turn-on bias (one draw at run start),
    then one white draw per sample. Within a fixed configuration the
    schedule is unconditional; enabling or disabling the optional
    components is a configuration change and lands in the resolved
    config hash (FR-15), which is the reproducibility anchor.
  \item \textbf{Geodetic reuse.} The altimeter calls the same core
    geodetic-altitude function the atmosphere and environment models
    use (Chapters~\ref{ch:harrispriester}, \ref{ch:environment}) --- no
    second conversion path; its accuracy budget
    (sub-millimetre in-domain) is inherited and already documented
    there.
  \item \textbf{Logged channels.} Measurements and flags land under
    \texttt{sensors.navfix.*} and \texttt{sensors.alt.*} with
    unit-suffixed names at the sensor rate.
\end{enumerate}

\section{Validation evidence}
\label{sec:radio:validation}

This chapter precedes its modules; the validation-evidence table with
machine-checked test identifiers is added by the implementing change.
The planned gates, mirroring Phase~6 exit criterion~6:
\begin{itemize}
  \item \textbf{Chi-square gates (exit criterion 6).} The three
    ensemble-mean statistics of
    equations~\eqref{eq:radio:chi2}--\eqref{eq:radio:bounds} over
    1{,}000 seeded draws fall inside their stated bounds, and rerunning
    the same seeds reproduces the statistics bit-for-bit.
  \item \textbf{Deterministic error terms.} With white terms zeroed, the
    altimeter output equals $h(\vv{r}) + b_h$ with $b_h$ the exact
    seeded stationary draw; the enabled Gauss--Markov component
    reproduces the recursion~\eqref{eq:radio:gm} value-for-value
    against an independent Python evaluation of the same stream.
  \item \textbf{Geodetic consistency.} On the reference ellipsoid
    ($h = 0$ states) and on a spherical body the noise-free altimeter
    returns $0$ and $\norm{\vv{r}} - a$ respectively to the
    Chapter~\ref{ch:harrispriester} conversion budget.
  \item \textbf{Gating.} Truth trajectories crossing $h_{\min}$ and
    $h_{\max}$ flip the altimeter flag exactly at the thresholds, with
    draws consumed identically on both sides.
  \item \textbf{Bit identity.} Two runs of the same seeded scenario
    produce byte-identical \texttt{sensors.navfix.*} and
    \texttt{sensors.alt.*} channels.
\end{itemize}

\section{References}
\label{sec:radio:references}

The fix-level abstraction of a navigation service and the additive
Gaussian error convention follow standard aided-navigation practice
(Titterton and Weston~\cite{titterton2004}); the Gauss--Markov machinery
is specified in Chapter~\ref{ch:sensors-imu} after
Maybeck~\cite{maybeck1979}; the geodetic conversion is
Bowring~\cite{bowring1976} as implemented in
Chapter~\ref{ch:harrispriester}. Full bibliographic details appear in
the bibliography.
